\documentclass[a4paper,11pt]{ltjsarticle}
\usepackage[haranoaji]{luatexja-preset}
\usepackage{amsmath,amssymb}
\usepackage[top=12mm,bottom=10mm,left=14mm,right=14mm]{geometry}
\usepackage{array}
\usepackage{tikz}
\usetikzlibrary{calc}
\pagestyle{empty}
\setlength{\parindent}{0pt}
\newlength{\qcol}\setlength{\qcol}{0.47\textwidth}
\newlength{\qgap}
\newlength{\anscolw}\setlength{\anscolw}{0.30\textwidth}
\newlength{\ansh}
\newcommand{\Q}[2]{\parbox[t]{\qcol}{\raggedright (#1)\hspace{0.6em}$\displaystyle #2$}}
\newcommand{\QR}[4]{\Q{#1}{#2}\hfill\Q{#3}{#4}\par\vspace{\qgap}}
\newcommand{\anscell}[1]{\rule[-\ansdrop]{0pt}{\ansh}(#1)}
\newlength{\ansdrop}

\begin{document}
{\large\bfseries 計算問題　22}\hfill 名前(\underline{\hspace{32mm}})\par\vspace{3mm}
■（1）〜（16）の計算をしなさい。（17）〜（21）は方程式を解きなさい。\par\vspace{3mm}
\setlength{\qgap}{5.4mm}
\QR{1}{\dfrac{9}{3}\div8+\dfrac{6}{3}}{2}{\dfrac{9}{3}(3x-18)-\dfrac{8}{6}(6x+54)}
\QR{3}{(-13)\times(-5)}{4}{-0.9-(-8.3)-1.6}
\QR{5}{(27a-24)\div3}{6}{(-27)\div(+3)}
\QR{7}{-9^2+(-3)^2-8}{8}{9(3x-2)-(8x-6)}
\QR{9}{(-16)+(-6)-(-13)}{10}{(3x-9)-(6x+8)}
\QR{11}{\dfrac{1}{9}(27x-72)+\dfrac{1}{6}(48x+18)}{12}{10\times\{3-(-8+1)\times2\}}
\QR{13}{17-(-3)\div\left(-\dfrac{3}{8}\right)}{14}{(9a+3)+(8a-6)}
\QR{15}{\dfrac{x-9}{3}\times24}{16}{9\div3\times(-24)}
\QR{17}{9x-3=-39}{18}{9(x+3)=10x+29}
\QR{19}{0.9x+0.3=-2.4}{20}{\dfrac{2}{3}x+1=\dfrac{1}{6}x-2}
\vspace{1mm}
\Q{21}{\dfrac{9x-3}{8}-\dfrac{x+6}{8}=-3}\par\vspace{3mm}
\setlength{\ansh}{9.5mm}\setlength{\ansdrop}{6mm}%
\begin{center}\begin{tabular}{|p{\anscolw}|p{\anscolw}|p{\anscolw}|}\hline
\anscell{1} & \anscell{2} & \anscell{3} \\\hline
\anscell{4} & \anscell{5} & \anscell{6} \\\hline
\anscell{7} & \anscell{8} & \anscell{9} \\\hline
\anscell{10} & \anscell{11} & \anscell{12} \\\hline
\anscell{13} & \anscell{14} & \anscell{15} \\\hline
\anscell{16} & \anscell{17} & \anscell{18} \\\hline
\anscell{19} & \anscell{20} & \anscell{21} \\\hline
\end{tabular}\end{center}

\end{document}
